% GENERATED FILE. Edit canonical lessons, then run npm run generate:books.
% canonical-source-hash: fnv1a64:3b9657ec10b4a61d
% canonical-lessons: LA-C453-ductor, LA-C453-gnarus, LA-C453-prosequor, LA-C453-reperio, LA-C453-vestigo

\chapter{Leader, Knowing, To escort, To find, To track}
\label{ch:la-leader-knowing-to-escort-to-find-to-track}
\hlchaptermodality{453}

\begin{chapteropening}
\textbf{By the end of this chapter:} \emph{I can talk about people who help you find the way: a leader, someone knowing the country, escorting, finding, and tracking.}

Complete the last lesson of chapter 453: I can talk about people who help you find the way: a leader, someone knowing the country, escorting, finding, and tracking.
\end{chapteropening}

\section[ductor, ductōris]{ductor, ductōris — a leader, one who leads the way; a commander}

\label{lesson:LA-C453-ductor}



\begin{quote}
Before the new one: say the Latin for a roundabout way, then the Latin for to wander about.
\end{quote}

\subsection*{You'll want to know: ductor, ductōris}
\textbf{ductor, ductōris} — \textquotedblleft{}a leader, one who leads the way; a commander\textquotedblright{}.

\begin{quote}
Ductor viātōrēs per montēs dūcit.
\end{quote}

\begin{quote}
\emph{The leader takes the travellers through the mountains.}
\end{quote}

Third declension, masculine: \emph{ductor, ductōris m.}, from \emph{dūcō} 'to lead'. \emph{Dux itineris} is the usual phrase for a guide.

\begin{cousinweb}
\emph{ductor} is built on \emph{dūcō} 'to lead' with the suffix \emph{-tor} 'one who' (EtymDB; Lewis \& Short). English \emph{conductor} comes from a compound of the same verb.
\end{cousinweb}

\begin{grammarlens}[title={What you've built}]
One more word for telling, explaining and describing at B1.
\end{grammarlens}

\subsection*{Your turn}
\begin{itemize}
\raggedright
  \item \emph{Say it:} \emph{ductor, ductōris}
  \item \emph{Say it:} \emph{ductor, ductōris}, once more
  \item \emph{Say it:} a sentence of your own with \emph{ductor, ductōris}
  \item \emph{Recall:} say the Latin for a roundabout way, then the Latin for to wander about, then say \emph{ductor, ductōris} again
\end{itemize}

\begin{tcolorbox}[breakable,colback=teal!4,colframe=teal!35!black,title={Before you move on}]
What does ductor, ductōris mean? (\textquotedblleft{}A leader, one who leads the way; a commander\textquotedblright{}.) Say it once more.
\end{tcolorbox}
\section[gnārus, gnāra, gnārum]{gnārus, gnāra, gnārum — knowing, familiar with, expert in}

\label{lesson:LA-C453-gnarus}



\begin{quote}
Before the new one: say the Latin for to wander about, then the Latin for a leader.
\end{quote}

\subsection*{You'll want to know: gnārus, gnāra, gnārum}
\textbf{gnārus, gnāra, gnārum} — \textquotedblleft{}knowing, familiar with, expert in\textquotedblright{}.

\begin{quote}
Pastor regiōnis gnārus est et viam scit.
\end{quote}

\begin{quote}
\emph{The shepherd knows the region well and knows the way.}
\end{quote}

An adjective of three endings, with a genitive: \emph{locōrum gnārus}, who knows the places. The opposite is \emph{ignārus}.

\begin{cousinweb}
\emph{gnārus} is inherited and comes from the same root as \emph{(g)nōscō} 'to get to know' and \emph{nārrō} 'to tell' (Lewis \& Short), from PIE \emph{ǵneh\textsubscript{3}-} 'to know'. English \emph{know} is a distant cousin.
\end{cousinweb}

\begin{grammarlens}[title={What you've built}]
One more word for telling, explaining and describing at B1.
\end{grammarlens}

\subsection*{Your turn}
\begin{itemize}
\raggedright
  \item \emph{Say it:} \emph{gnārus, gnāra, gnārum}
  \item \emph{Say it:} \emph{gnārus, gnāra, gnārum}, once more
  \item \emph{Say it:} a sentence of your own with \emph{gnārus, gnāra, gnārum}
  \item \emph{Recall:} say the Latin for to wander about, then the Latin for a leader, then say \emph{gnārus, gnāra, gnārum} again
\end{itemize}

\begin{tcolorbox}[breakable,colback=teal!4,colframe=teal!35!black,title={Before you move on}]
What does gnārus, gnāra, gnārum mean? (\textquotedblleft{}Knowing, familiar with, expert in\textquotedblright{}.) Say it once more.
\end{tcolorbox}
\section[prōsequor, prōsequī]{prōsequor, prōsequī — to escort, to see off, to accompany; to pursue}

\label{lesson:LA-C453-prosequor}



\begin{quote}
Before the new one: say the Latin for a leader, then the Latin for knowing.
\end{quote}

\subsection*{You'll want to know: prōsequor, prōsequī}
\textbf{prōsequor, prōsequī} — \textquotedblleft{}to escort, to see off, to accompany; to pursue\textquotedblright{}.

\begin{quote}
Amīcī viātōrem ad portam urbis prōsequuntur.
\end{quote}

\begin{quote}
\emph{Friends escort the traveller as far as the city gate.}
\end{quote}

Deponent: \emph{prōsequor, prōsequī, prōsecūtus sum}. Escorting a departing friend part of the way was a duty. English \emph{prosecute} and \emph{pursue} come from it.

\begin{cousinweb}
\emph{prōsequor} is built from \emph{prō-} 'forward, onward' and \emph{sequor} 'to follow' (EtymDB; Lewis \& Short). English \emph{prosecute} comes from its participle, and \emph{pursue} from it through French.
\end{cousinweb}

\begin{grammarlens}[title={What you've built}]
One more word for telling, explaining and describing at B1.
\end{grammarlens}

\subsection*{Your turn}
\begin{itemize}
\raggedright
  \item \emph{Say it:} \emph{prōsequor, prōsequī}
  \item \emph{Say it:} \emph{prōsequor, prōsequī}, once more
  \item \emph{Say it:} a sentence of your own with \emph{prōsequor, prōsequī}
  \item \emph{Recall:} say the Latin for a leader, then the Latin for knowing, then say \emph{prōsequor, prōsequī} again
\end{itemize}

\begin{tcolorbox}[breakable,colback=teal!4,colframe=teal!35!black,title={Before you move on}]
What does prōsequor, prōsequī mean? (\textquotedblleft{}To escort, to see off, to accompany; to pursue\textquotedblright{}.) Say it once more.
\end{tcolorbox}
\section[reperiō, reperīre]{reperiō, reperīre — to find, to find out, to discover}

\label{lesson:LA-C453-reperio}



\begin{quote}
Before the new one: say the Latin for knowing, then the Latin for to escort.
\end{quote}

\subsection*{You'll want to know: reperiō, reperīre}
\textbf{reperiō, reperīre} — \textquotedblleft{}to find, to find out, to discover\textquotedblright{}.

\begin{quote}
Tandem viātōrēs caupōnam reperiunt.
\end{quote}

\begin{quote}
\emph{At last the travellers find an inn.}
\end{quote}

\emph{reperiō, reperīre, repperī, repertum}. English \emph{repertory} comes from \emph{repertōrium}, a list of things found.

\begin{cousinweb}
\emph{reperiō} is built from \emph{re-} and \emph{pariō} 'to bring forth, to get' (Lewis \& Short): to get again, to find. English \emph{repertory} and \emph{repertoire} come from \emph{repertōrium}, built on it.
\end{cousinweb}

\begin{grammarlens}[title={What you've built}]
One more word for telling, explaining and describing at B1.
\end{grammarlens}

\subsection*{Your turn}
\begin{itemize}
\raggedright
  \item \emph{Say it:} \emph{reperiō, reperīre}
  \item \emph{Say it:} \emph{reperiō, reperīre}, once more
  \item \emph{Say it:} a sentence of your own with \emph{reperiō, reperīre}
  \item \emph{Recall:} say the Latin for knowing, then the Latin for to escort, then say \emph{reperiō, reperīre} again
\end{itemize}

\begin{tcolorbox}[breakable,colback=teal!4,colframe=teal!35!black,title={Before you move on}]
What does reperiō, reperīre mean? (\textquotedblleft{}To find, to find out, to discover\textquotedblright{}.) Say it once more.
\end{tcolorbox}
\section[vestīgō, vestīgāre]{vestīgō, vestīgāre — to track, to trace; to search out}

\label{lesson:LA-C453-vestigo}



\begin{quote}
Before the new one: say the Latin for to escort, then the Latin for to find.
\end{quote}

\subsection*{You'll want to know: vestīgō, vestīgāre}
\textbf{vestīgō, vestīgāre} — \textquotedblleft{}to track, to trace; to search out\textquotedblright{}.

\begin{quote}
Canis vestīgia lupī in nive vestīgat.
\end{quote}

\begin{quote}
\emph{The dog tracks the wolf's footprints in the snow.}
\end{quote}

\emph{vestīgō, vestīgāre, vestīgāvī, vestīgātum}. English \emph{investigate} comes from \emph{investīgāre}.

\begin{cousinweb}
\emph{vestīgō} goes with \emph{vestīgium} 'footprint, track'; Lewis \& Short call the origin doubtful, so which came first is not settled. English \emph{investigate} comes from the compound \emph{investīgāre}.
\end{cousinweb}

\begin{grammarlens}[title={What you've built}]
One more word for telling, explaining and describing at B1.
\end{grammarlens}

\subsection*{Your turn}
\begin{itemize}
\raggedright
  \item \emph{Say it:} \emph{vestīgō, vestīgāre}
  \item \emph{Say it:} \emph{vestīgō, vestīgāre}, once more
  \item \emph{Say it:} a sentence of your own with \emph{vestīgō, vestīgāre}
  \item \emph{Recall:} say the Latin for to escort, then the Latin for to find, then say \emph{vestīgō, vestīgāre} again
\end{itemize}

\begin{tcolorbox}[breakable,colback=teal!4,colframe=teal!35!black,title={Before you move on}]
What does vestīgō, vestīgāre mean? (\textquotedblleft{}To track, to trace; to search out\textquotedblright{}.) Say it once more.
\end{tcolorbox}
